\documentclass[10pt,a4paper]{amsart}
\usepackage{amsmath,amssymb,amsthm,graphicx}
\usepackage{verbatim}
\newenvironment{reference}{\comment}{\endcomment}
%\newenvironment{reference}{}{}
\newenvironment{slogan}{\comment}{\endcomment}
\newenvironment{history}{\comment}{\endcomment}

% For commutative diagrams we use Xy-pic
\usepackage[all]{xy}

% We use 2cell for 2-commutative diagrams.
\xyoption{2cell}
% Theorem environments.
%
\theoremstyle{plain}
\newtheorem{theorem}[subsection]{Theorem}
\newtheorem{proposition}[subsection]{Proposition}
\newtheorem{lemma}[subsection]{Lemma}

\theoremstyle{definition}
\newtheorem{definition}[subsection]{Definition}
\newtheorem{example}[subsection]{Example}
\newtheorem{exercise}[subsection]{Exercise}
\newtheorem{situation}[subsection]{Situation}

\theoremstyle{remark}
\newtheorem{remark}[subsection]{Remark}
\newtheorem{remarks}[subsection]{Remarks}

\numberwithin{equation}{subsection}

% Macros
%
\def\lim{\mathop{\mathrm{lim}}\nolimits}
\def\colim{\mathop{\mathrm{colim}}\nolimits}
\def\Spec{\mathop{\mathrm{Spec}}}
\def\Hom{\mathop{\mathrm{Hom}}\nolimits}
\def\Ext{\mathop{\mathrm{Ext}}\nolimits}
\def\SheafHom{\mathop{\mathcal{H}\!\mathit{om}}\nolimits}
\def\SheafExt{\mathop{\mathcal{E}\!\mathit{xt}}\nolimits}
\def\Sch{\mathit{Sch}}
\def\Mor{\mathop{\mathrm{Mor}}\nolimits}
\def\Ob{\mathop{\mathrm{Ob}}\nolimits}
\def\Sh{\mathop{\mathit{Sh}}\nolimits}
\def\NL{\mathop{N\!L}\nolimits}
\def\CH{\mathop{\mathrm{CH}}\nolimits}
\def\proetale{{pro\text{-}\acute{e}tale}}
\def\etale{{\acute{e}tale}}
\def\QCoh{\mathit{QCoh}}
\def\Ker{\mathop{\mathrm{Ker}}}
\def\Im{\mathop{\mathrm{Im}}}
\def\Coker{\mathop{\mathrm{Coker}}}
\def\Coim{\mathop{\mathrm{Coim}}}

% Boxtimes
%
\DeclareMathSymbol{\boxtimes}{\mathbin}{AMSa}{"02}

%
% Macros for moduli stacks/spaces
%
\def\QCohstack{\mathcal{QC}\!\mathit{oh}}
\def\Cohstack{\mathcal{C}\!\mathit{oh}}
\def\Spacesstack{\mathcal{S}\!\mathit{paces}}
\def\Quotfunctor{\mathrm{Quot}}
\def\Hilbfunctor{\mathrm{Hilb}}
\def\Curvesstack{\mathcal{C}\!\mathit{urves}}
\def\Polarizedstack{\mathcal{P}\!\mathit{olarized}}
\def\Complexesstack{\mathcal{C}\!\mathit{omplexes}}
% \Pic is the operator that assigns to X its picard group, usage \Pic(X)
% \Picardstack_{X/B} denotes the Picard stack of X over B
% \Picardfunctor_{X/B} denotes the Picard functor of X over B
\def\Pic{\mathop{\mathrm{Pic}}\nolimits}
\def\Picardstack{\mathcal{P}\!\mathit{ic}}
\def\Picardfunctor{\mathrm{Pic}}
\def\Deformationcategory{\mathcal{D}\!\mathit{ef}}
\usepackage{cleveref}
\crefname{originallemma}{original lemma}{original lemmas}
\Crefname{originallemma}{Original lemma}{Original lemmas}
\makeatletter
\newlabel{lemma-cocontinuous-morphism-topoi}{{00XO}{source}{}{sourcepointer-1}{}}
\newlabel{lemma-cocontinuous-morphism-topoi@cref}{{[originallemma][0][]00XO}{[1][0][]source}{}{}{}}
\makeatother
\crefname{originaldefinition}{original definition}{original definitions}
\Crefname{originaldefinition}{Original definition}{Original definitions}
\makeatletter
\newlabel{definition-morphism-sites}{{00X1}{source}{}{sourcepointer-2}{}}
\newlabel{definition-morphism-sites@cref}{{[originaldefinition][0][]00X1}{[1][0][]source}{}{}{}}
\makeatother
\title{A continuous right adjoint of a cocontinuous functor defines the same morphism of topoi}
\author{The Stacks Project Authors}
\date{stacks-project:sites.tex:local-snapshot-20261012; GNU FDL 1.2 or later}
\begin{document}
\maketitle
\noindent\textit{Editorial note.} All delivered mathematical text is the author's own, copied from the pinned Stacks Project chapter file sites-cohomology.tex; no statement or proof step was rewritten, reordered or removed.\par\medskip
\noindent\textit{Editorial note.} The proof is the author's complete argument; every step the author writes out is inside the counted span and no closing sentence is a placeholder.\par\medskip
\noindent\textit{Editorial note.} Every labelled reference occurring inside the counted spans is carried in assembly.external\_references as a non-load-bearing navigation pointer (same-chapter ones at their definition site, cross-chapter ones at the invocation inside the counted span), so no load-bearing mathematics is replaced by a stub. No bibliography citation occurs inside the counted spans.\par\medskip
\noindent\textit{Editorial note.} The checked chapter list for same-name collisions is the set of stacks chapters staged in this repository: chow, cohomology, divisors, etale-cohomology, flat, more-algebra, more-morphisms, morphisms, spaces-more-morphisms, varieties together with this card's chapter sites-cohomology.\par\medskip
\noindent\textit{Editorial note.} The preamble is the chapter's own retained minimal preamble (sources/stacks-sites-cohomology/retained-preamble.tex), a verbatim excerpt of the native preamble: the Xy-pic setup and the native theorem declarations with the mathematical macros.\par\medskip
\noindent\textit{Allowed standard prerequisites (editorial).}\begin{itemize}
\item The morphism of topoi associated to a cocontinuous functor and the description of its functors. Same chapter, Lemma lemma-have-functor-other-way, recorded as the context of this card. Supplies the morphism of topoi and the identifications used in the proof.
\item The definition of a morphism of sites. Same chapter, Definition definition-morphism-sites. Supplies the exactness condition that has to be checked for the right adjoint.
\end{itemize}\medskip
\ifdefined\hypertarget\hypertarget{sourcepointer-1}{}\fi
\noindent\textit{Source reference.} Original lemma 00XO: sources/stacks-sites/pinned-original-source.tex, source lines 5063--5063. Navigation pointer invoked at proof line 5063; the counted span contains the pointer only and the cited result is not reproduced inside it, so the pointer is registered as a non-load-bearing navigation stub in the established form for Stacks material. No bibliography citation is involved.\par\medskip
\ifdefined\hypertarget\hypertarget{sourcepointer-2}{}\fi
\noindent\textit{Source reference.} Original definition 00X1: sources/stacks-sites/pinned-original-source.tex, source lines 5109--5109. Navigation pointer invoked at proof line 5109; the counted span contains the pointer only and the cited result is not reproduced inside it, so the pointer is registered as a non-load-bearing navigation stub in the established form for Stacks material. No bibliography citation is involved.\par\medskip
\begin{lemma}
\label{lemma-have-functor-other-way}
Let $\mathcal{C}$ and $\mathcal{D}$ be sites. Let
$u : \mathcal{C} \to \mathcal{D}$, and $v : \mathcal{D} \to \mathcal{C}$
be functors. Assume that $u$ is cocontinuous
and that $v$ is a right adjoint to $u$.
Let $g : \Sh(\mathcal{C}) \to \Sh(\mathcal{D})$ be
the morphism of topoi associated to $u$, see
Lemma \ref{lemma-cocontinuous-morphism-topoi}.
Then
\begin{enumerate}
\item for a sheaf $\mathcal{F}$ on $\mathcal{C}$ the sheaf
$g_*\mathcal{F}$ is equal to the presheaf $v^p\mathcal{F}$, in other words,
$(g_*\mathcal{F})(V) = \mathcal{F}(v(V))$, and
\item for a sheaf $\mathcal{G}$ on $\mathcal{D}$ we have
$g^{-1}\mathcal{G} = (v_p\mathcal{G})^\#$.
\end{enumerate}
\end{lemma}

\begin{lemma}
\label{lemma-have-functor-other-way-morphism}
Notation and assumptions as in Lemma \ref{lemma-have-functor-other-way}.
If in addition $v$ is continuous then $v$ defines a morphism of sites
$f : \mathcal{C} \to \mathcal{D}$ whose associated morphism
of topoi is equal to $g$.
\end{lemma}

\begin{proof}
We will use the results of Lemma \ref{lemma-have-functor-other-way}
without further mention. To prove that $v$ defines a morphism of sites
$f$ as in the statement of the lemma,
we have to show that $v_s$ is an exact functor (see
Definition \ref{definition-morphism-sites}). Since
$v_s\mathcal{G} = (v_p\mathcal{G})^\# = g^{-1}\mathcal{G}$
this follows from the fact that $g$ is a morphism of topoi.
Then we see that $f^{-1} = v_s = g^{-1}$ and we find that $f = g$
as morphisms of topoi.
\end{proof}

\end{document}
